\section{{\fontsize{12pt}{14pt}\selectfont REFERENCIAL TEÓRICO}}
\vspace{-0.3cm}
\novoBlocoTexto{
    Este capítulo reúne os conceitos que fundamentam este trabalho e que apoiam a compreensão do assistente proposto e das escolhas feitas nos capítulos seguintes. Ele está organizado em seis seções. A seção~\ref{sec:duvidas-discentes} trata das dúvidas discentes, da busca de ajuda acadêmica e do uso dessas dúvidas como subsídio à prática docente. A seção~\ref{sec:chatbots} apresenta os \textit{chatbots} e o seu papel na educação, tanto no apoio ao estudante quanto como fonte de informação para o professor. As seções~\ref{sec:llm} e~\ref{sec:rag} descrevem as LLM e o RAG, tecnologias em que se baseia o assistente proposto. A seção~\ref{sec:classificacao-duvidas-materiais} aborda a classificação de dúvidas a partir do material didático do professor, com as representações vetoriais, a classificação textual por LLM e as métricas de avaliação. Por fim, a seção~\ref{sec:privacidade} discute a privacidade e a proteção dos dados dos estudantes.
}
\subsection{{\fontsize{12pt}{14pt}\selectfont Dúvidas discentes}}
\label{sec:duvidas-discentes}
\novoBlocoTexto{
    As dúvidas fazem parte do processo de aprendizagem discente, podendo surgir à medida que os estudantes entram em contato com novos conhecimentos e buscam compreendê-los.
}
\subsubsection{{\fontsize{12pt}{14pt}\selectfont Busca de ajuda acadêmica}}
\novoBlocoTexto{
    Diante do cenário de dúvidas discentes, a busca de ajuda acadêmica apresenta-se como uma estratégia de autorregulação que contribui para este processo de aprendizagem \parencite{ryan2001avoiding}. Porém, diferentemente de outras estratégias autorregulatórias, a busca de ajuda acadêmica apresenta uma dimensão interpessoal, uma vez que envolve necessariamente a comunicação ou a interação com outras pessoas \parencite{ryan1997should}.
}
\novoBlocoTexto{
    Historicamente, a busca de ajuda foi associada à dependência e considerada uma prática degradante, frequentemente evitada \parencite{nelsonlegall1985help}. Essa percepção também se manifesta no contexto universitário, no qual muitos estudantes deixam de solicitar o auxílio necessário por considerarem essa atitude constrangedora ou um reconhecimento de fracasso \parencite{karabenick2011understanding}. Entre os fatores que dificultam essa iniciativa, destacam-se o receio de julgamentos por parte de professores e colegas \parencite{ryan1997should}, a indisponibilidade de pessoas que possam oferecer auxílio e o tempo de espera por uma resposta \parencite{ryan2001avoiding}.
}
\novoBlocoTexto{
    Além disso, pesquisas sobre a busca de ajuda em ambientes escolares indicam que estudantes com menor desempenho acadêmico podem apresentar maior resistência em solicitar auxílio, mesmo quando necessitam de apoio \parencite{ryan2001avoiding}. Nesse sentido, a ausência de perguntas não significa necessariamente a inexistência de dificuldades de aprendizagem, pois o estudante pode não compreender suficientemente o conteúdo para formular uma dúvida ou pode sentir constrangimento em manifestá-la \parencite{karabenick2011understanding}. Como consequência, as dificuldades enfrentadas pelos estudantes podem continuar desconhecidas pelos docentes, o que os limita na identificação das necessidades de aprendizagem dos alunos.
}
\subsubsection{{\fontsize{12pt}{14pt}\selectfont Subsídios à prática docente}}
\novoBlocoTexto{
    As dúvidas manifestadas pelos estudantes podem representar necessidades individuais de aprendizagem, como também, informações relevantes para a prática docente. \textcite{freire2013pedagogia} ressaltam que o professor também aprende com aqueles a quem ensina, na medida em que as inquietações e os questionamentos dos estudantes o desafiam a reconsiderar seus próprios conhecimentos. Essa perspectiva relaciona-se ao conceito de avaliação formativa, no qual as evidências sobre a aprendizagem são obtidas, interpretadas e utilizadas para orientar decisões sobre os próximos passos do processo de ensino e aprendizagem \parencite{black2009developing}. 
}
\novoBlocoTexto{
    Em um dos casos relatados por \textcite{angelo1993classroom}, uma professora identificou que sua abordagem estava baseada na forma como havia aprendido o conteúdo durante a pós-graduação, sem considerar suficientemente as necessidades de seus estudantes, o que a levou a modificar as aulas seguintes. Os autores reforçam a necessidade de obter informações continuamente sobre a aprendizagem, considerando que a identificação tardia das dificuldades pode limitar as possibilidades de intervenção pedagógica \parencite{angelo1993classroom}.
}
\novoBlocoTexto{
    Entre as estratégias destinadas à identificação de dificuldades discentes, \textcite{angelo1993classroom} apresentam a técnica denominada \textit{muddiest point}, que consiste em solicitar aos estudantes que indiquem os aspectos menos compreendidos de uma aula ou atividade. As respostas permitem identificar conteúdos que demandam esclarecimentos adicionais e orientar a distribuição do tempo dedicado a determinados tópicos. Para isso, os autores propõem o agrupamento das respostas por assuntos relacionados e a identificação da frequência de ocorrência de cada grupo \parencite{angelo1993classroom}. 
}
\novoBlocoTexto{
    Entretanto, a disponibilidade de informações sobre o andamento da aprendizagem discente não assegura, por si só, sua utilização no aperfeiçoamento do ensino. \textcite{wise2019teaching} evidenciam a complexidade do processo de tomada de decisão docente fundamentada em dados, envolvendo sua interpretação, a definição de ações e a avaliação dos resultados obtidos. Dessa forma, as dúvidas discentes podem carregar subsídios importantes para a prática pedagógica, desde que sejam analisadas e contextualizadas adequadamente. Sua organização pode contribuir para a identificação de conteúdos que exigem maior atenção, enquanto a definição das intervenções didáticas permanece sob responsabilidade do professor.
}
\subsection{{\fontsize{12pt}{14pt}\selectfont\textit{Chatbots}}}
\label{sec:chatbots}
\novoBlocoTexto{
    \textit{Chatbots} são sistemas computacionais desenvolvidos para interagir com usuários por meio da linguagem natural. Também são conhecidos como agentes conversacionais, agentes interativos, \textit{bots} inteligentes ou assistentes digitais \parencite{adamopoulou2020chatbots}.}    
\subsubsection{{\fontsize{12pt}{14pt}\selectfont Função dos \textit{chatbots} no contexto educacional}}
\novoBlocoTexto{
    Na perspectiva educacional, os \textit{chatbots} podem oferecer suporte pedagógico de forma escalável, respondendo dúvidas conceituais e operacionais, oferecendo \textit{feedback} formativo e orientando trajetórias de estudo \parencite{guimaraes2026agentes}.
}
\novoBlocoTexto{
    \textcite{oliveira2026agentes} observam que os modelos LLM ampliaram as possibilidades de personalização e de interação em linguagem natural. Em comparação com sistemas conversacionais baseados em regras predefinidas, os agentes fundamentados em LLM apresentam maior flexibilidade na geração de respostas, possibilitando interações mais adaptáveis às necessidades dos usuários. Nesse cenário, os \textit{chatbots} passaram a ser o tipo de agente educacional mais estudado, presentes em 68,96\% dos 58 estudos revisados por \textcite{elnaffar2026teaching}. 
}
\novoBlocoTexto{
    \textcite{kuhail2023interacting} identificam quatro papéis desses agentes junto ao estudante: o de agente de ensino, que faz o papel do professor ao apresentar instruções, dar exemplos, fazer perguntas e oferecer \textit{feedback} imediato; o de agente par, que atua como colega de estudo; o de agente ensinável, que é ensinado pelo próprio aluno; e o de agente motivacional, que acompanha o estudante e incentiva a aprendizagem. Entre esses papéis, o mais frequente é o de agente de ensino, presente em 55,55\% dos \textit{chatbots} analisados pelos autores \parencite{kuhail2023interacting}.
}
\subsubsection{{\fontsize{12pt}{14pt}\selectfont \textit{Chatbots} no processo de aprendizagem discente}}
\novoBlocoTexto{
    Uma das principais contribuições desses assistentes no processo de aprendizagem discente é oferecer apoio no momento em que a dúvida surge. Na revisão de \textcite{elnaffar2026teaching}, os principais benefícios relatados foram o \textit{feedback} personalizado e o acesso a qualquer hora, que apoiam o estudo no próprio ritmo e reduzem a frustração do aluno. 
}
\novoBlocoTexto{
    Outro papel relevante diz respeito à exposição do estudante. Ao conversar com o assistente, o estudante diminui o receio de ser julgado pelas perguntas que faz \parencite{ma2024integrating}. \textcite{liffiton2023codehelp} relatam que alunos que pareciam desconfortáveis, envergonhados ou relutantes em pedir ajuda ao professor ou ao monitor não tiveram receio de perguntar ao assistente. A própria conversa também contribui para a aprendizagem, pois muitas vezes o estudante não conhece o bastante sobre a própria dúvida para formulá-la adequadamente ou mesmo classificá-la, e o diálogo com o tutor pode ajudá-lo a descobrir qual é, de fato, a sua pergunta \parencite{ma2024integrating}.
}
\novoBlocoTexto{
    Quanto aos efeitos sobre a aprendizagem, uma meta-análise recente aponta resultados favoráveis. \textcite{deng2023metaanalysis} encontraram um efeito global de 0,789 em 32 estudos, com efeito maior quando o \textit{chatbot} atuava como assistente de ensino \parencite{deng2023metaanalysis}. Em um experimento conduzido por \textcite{kestin2025ai}, estudantes que utilizaram um tutor baseado em inteligência artificial (IA, do inglês, \textit{Artificial Intelligence}) apresentaram maior aprendizagem em menor tempo, quando comparados àqueles que participaram de atividades de aprendizagem ativa em sala de aula. Entretanto, os autores ressaltam que o tutor não deve substituir as aulas presenciais, mas preparar os estudantes para aproveitá-las melhor \parencite{kestin2025ai}.
}
\novoBlocoTexto{
    É importante ressaltar que o uso desses assistentes como ferramenta educacional também envolve riscos. Entre os principais desafios identificados por \textcite{elnaffar2026teaching} destacam-se a dependência excessiva, acompanhada de aprendizagem superficial, relatada em 65,52\% dos estudos revisados, seguida das respostas incorretas, inconsistentes ou incoerentes (32,76\%). \textcite{deng2023metaanalysis} registraram ainda em sua meta-análise resultados negativos quanto ao pensamento crítico, ao engajamento e à motivação dos alunos. Diante disso, essas ferramentas não devem substituir o professor, mas sim ser utilizadas de forma complementar ao ensino.
}
\subsubsection{{\fontsize{12pt}{14pt}\selectfont Interações discentes com \textit{chatbots} como subsídios à prática docente}}
\novoBlocoTexto{
    Conforme a descrição dos assistentes na literatura, nota-se que eles se dirigem majoritariamente ao estudante. \textcite{adamopoulou2020chatbots} chegam a mencionar que os \textit{chatbots} reúnem informações ao longo da disciplina que ajudam a melhorar a aprendizagem e o ensino, mas esse uso é pouco explorado. \textcite{alves2024chatbots} encontraram 14 estudos com estudantes e apenas 3 com professores, e apontam a falta de estudos centrados na perspectiva docente como uma lacuna da área \parencite{alves2024chatbots}.
}
\novoBlocoTexto{
    \textcite{gamito2019botanalyzer} observa que, ao oferecer aos alunos um ambiente disponível a qualquer hora para tirar dúvidas, o \textit{chatbot} alivia a carga de trabalho do professor, mas faz com que ele perca contato com os estudantes, o que gera uma lacuna entre as dúvidas que eles têm e o conhecimento que o professor tem delas. Segundo o autor, a única forma de retorno disponível ao docente costuma ser o registro bruto das interações, o que dificulta identificar as dúvidas comuns à turma e os assuntos mais difíceis de serem compreendidos \parencite{gamito2019botanalyzer}. O próprio autor, inclusive, desenvolveu uma ferramenta, o \textit{BotAnalyzer}, que visava facilitar ao professor a visualização das dúvidas e dos problemas que os alunos tiveram durante os estudos, para que ele pudesse se preparar com antecedência para as aulas sem precisar ler um extenso arquivo de registros \parencite{gamito2019botanalyzer}. 
}
\novoBlocoTexto{
    Outras ferramentas também dão ao professor acesso às perguntas dos alunos. O \textit{CodeHelp} é um assistente de programação baseado em LLM, com barreiras de proteção (\textit{guardrails}) que o impedem de entregar as soluções prontas, e foi utilizado durante 12 semanas por 52 estudantes de uma disciplina do primeiro ano de computação e ciência de dados \parencite{liffiton2023codehelp}. Pela leitura das consultas feitas ao assistente, o professor obteve uma visão mais direta do raciocínio dos alunos do que a oferecida pelos trabalhos entregues \parencite{liffiton2023codehelp}. O \textit{MerryQuery}, por sua vez, é um assistente educacional baseado em LLM e RAG que responde aos estudantes com base nos materiais da disciplina, indica as fontes de cada resposta e permite ao docente acompanhar o histórico das conversas \parencite{tabarsi2025merryquery}. Nas três ferramentas, porém, as perguntas não chegam ao professor organizadas pelos conteúdos da disciplina: o \textit{BotAnalyzer}, citado no parágrafo anterior, mostra as palavras mais frequentes nas perguntas, em uma nuvem de palavras \parencite{gamito2019botanalyzer}, enquanto o \textit{CodeHelp} e o \textit{MerryQuery} as apresentam como uma lista. É nessa lacuna que se insere o assistente proposto neste trabalho, que mantém a função de apoio ao estudante, mas tem como objetivo principal levar ao professor as dúvidas dos alunos, organizadas segundo os tópicos extraídos de seu próprio plano de ensino.
}
\subsection{{\fontsize{12pt}{14pt}\selectfont Modelos LLM}}
\label{sec:llm}
\novoBlocoTexto{
    Os modelos LLM são modelos de inteligência artificial geralmente fundamentados na arquitetura \textit{Transformer}, caracterizados pela utilização de redes neurais com bilhões de parâmetros. Esses modelos são pré-treinados em grandes volumes de dados textuais, por meio dos quais aprendem padrões e relações linguísticas, podendo posteriormente ser adaptados para diferentes aplicações de processamento de linguagem natural, como geração de textos, respostas a perguntas e classificação de informações \parencite{zhao2026survey}.
}
\novoBlocoTexto{
    Com o GPT-3, modelo de linguagem com 175 bilhões de parâmetros, \textcite{brown2020language} mostraram que uma LLM suficientemente grande pode realizar tarefas novas sem nenhum treinamento adicional, apenas a partir do texto de entrada. No modo \textit{few-shot}, expressão que pode ser traduzida como ``com poucos exemplos'', o modelo recebe, no próprio texto de entrada, a descrição da tarefa e alguns exemplos já resolvidos, como frases em inglês acompanhadas de suas traduções para o francês, seguidos de um novo caso a resolver, sem que seus parâmetros sejam atualizados. No modo \textit{zero-shot}, ou ``sem exemplos'', o modelo recebe apenas a descrição da tarefa em linguagem natural \parencite{brown2020language}. Esse segundo modo é retomado na seção~\ref{sec:classificacao-textual-llm}.
}
\novoBlocoTexto{
    Essa flexibilidade é um dos fatores que explica a variedade de aplicações dos modelos LLM. Na educação, além de apoiarem os estudantes, eles podem ajudar o professor a personalizar o ensino, a elaborar planos de aula e atividades e a avaliar trabalhos, destacando seus pontos fortes e fracos \parencite{kasneci2023chatgpt}.
}
\novoBlocoTexto{
    Contudo, um dos principais riscos associados ao uso de LLM é a ocorrência de alucinações, caracterizadas pela geração de informações incorretas, inconsistentes ou sem fundamentação factual, mesmo quando apresentadas de maneira coerente e convincente. Segundo \textcite{anhhoang2025survey}, esse fenômeno pode decorrer tanto de problemas na formulação das instruções fornecidas ao modelo quanto de limitações inerentes ao seu funcionamento, relacionadas aos dados utilizados no treinamento e ao processo de geração das respostas. Mesmo os criadores do GPT-3 reconhecem que o modelo às vezes perde a coerência em textos longos e contradiz a si mesmo \parencite{brown2020language}.
}
\novoBlocoTexto{
    Com o objetivo de reduzir a alucinação por meio de embasamento da resposta em fontes verificáveis, surgiu o RAG, que incorpora à LLM conhecimento de bases externas \parencite{gao2023retrieval,lewis2020retrieval}. Porém, é importante ressaltar que a utilização dessa estratégia não elimina a possibilidade de alucinações, embora possa reduzir sua ocorrência. Assim, a literatura recomenda a supervisão humana para acompanhar e avaliar os resultados dessas soluções. \textcite{kasneci2023chatgpt} a consideram indispensável para o uso benéfico das LLM na educação.
}
\subsection{{\fontsize{12pt}{14pt}\selectfont RAG}}
\label{sec:rag}
\novoBlocoTexto{
    O RAG se consolidou como forma de reduzir a alucinação e de incorporar conhecimento atualizado e específico de uma área diante da popularização dos modelos LLM \parencite{gao2023retrieval}. 
    \textcite{lewis2020retrieval} propuseram o RAG para enfrentar limitações dos modelos de linguagem que dependem exclusivamente do conhecimento armazenado em seus parâmetros, como a dificuldade de atualizar esse conhecimento, identificar a origem das informações e evitar a geração de conteúdos incorretos. Essa estratégia combina a memória interna da LLM, chamada paramétrica, com uma memória externa, não paramétrica, formada por documentos acessados por um mecanismo de recuperação \parencite{lewis2020retrieval}. 
}
\novoBlocoTexto{
    Com RAG, primeiro um mecanismo de recuperação busca informações relevantes em uma base de conhecimento externa e depois as fornece à LLM para orientar a geração da resposta. Esse processo contribui para a precisão, a atualidade e a transparência das informações, característica especialmente importante na educação, em que as informações precisam ser verificáveis \parencite{li2025retrieval}. Outra vantagem é que, como o conhecimento fica em uma base externa, ele pode ser atualizado apenas substituindo essa base, sem treinar o modelo novamente \parencite{lewis2020retrieval}. Com isso, o RAG também se destaca por se adaptar a conteúdos que mudam com frequência.
}
\novoBlocoTexto{
    Na educação, o RAG vem sendo aplicado sobretudo em sistemas de perguntas e respostas, aprendizagem personalizada e tutoria, com ganhos de precisão, relevância e personalização das respostas \parencite{arthana2026systematic}. Nessas aplicações, a base de conhecimento é formada pelos próprios materiais didáticos, como livros, planos de ensino e listas de exercícios, que são convertidos em texto antes da indexação \parencite{li2025retrieval}. \textcite{nemeth2025exploring} destacam que a principal vantagem do RAG no ensino superior é justamente o controle do professor sobre a bibliografia usada pelo modelo, que passa a responder com conhecimento fundamentado em materiais selecionados pelo docente. Em uma disciplina introdutória de programação, \textcite{jacobs2024leveraging} usaram as transcrições das aulas como base do RAG, com o objetivo de que a LLM evitasse alucinações e empregasse os mesmos termos técnicos usados pelo professor.
}
\novoBlocoTexto{
    Apesar de suas potencialidades, o RAG apresenta limitações que precisam ser consideradas. A recuperação pode trazer trechos irrelevantes ou deixar de trazer informações essenciais, e o modelo ainda pode gerar conteúdo não sustentado pelo contexto recuperado \parencite{gao2023retrieval}. \textcite{swacha2025retrieval} acrescentam o custo computacional, já que cada pergunta exige uma etapa de recuperação e outra de geração. Essas limitações reforçam a necessidade de supervisão humana na utilização de sistemas educacionais baseados em RAG, a fim de identificar e corrigir possíveis inconsistências.
}
\subsection{{\fontsize{12pt}{14pt}\selectfont Classificação de dúvidas baseada em material didático docente}}
\label{sec:classificacao-duvidas-materiais}
\novoBlocoTexto{
    Por meio da organização de dúvidas por classificação baseada em tópicos ou seções do próprio plano de ensino docente, o professor pode identificar áreas em que os alunos mais pedem ajuda. O docente, então, pode ajustar seus métodos e intervenções, e o tipo de pergunta pode revelar lacunas no currículo ou no material didático \parencite{sheese2023patterns}. O obstáculo é o trabalho: categorizar manualmente todas as perguntas de uma disciplina é uma tarefa tediosa. Por isso, para tal tarefa, a literatura recorre a técnicas automáticas, como a classificação, em que cada texto é atribuído a uma categoria definida previamente.
}
\subsubsection{{\fontsize{12pt}{14pt}\selectfont Representações vetoriais e recuperação semântica}}
\novoBlocoTexto{
    Representações vetoriais e recuperação semântica são métodos para realizar comparação computacional entre textos. A forma mais tradicional de representação de vetores é a contagem de palavras, como a Frequência do Termo–Frequência Inversa do Documento (TF-IDF, do inglês \textit{Term Frequency–Inverse Document Frequency}), que combina a frequência de cada termo no texto com a sua raridade no conjunto, de modo que palavras muito comuns, como “é” ou “que”, pesem menos \parencite{gao2021automatically}. Essa representação, conhecida como “saco de palavras”, não captura diretamente as relações semânticas e contextuais entre as palavras \parencite{grootendorst2022bertopic}. Como alternativa a esse método, consideram-se os \textit{embeddings}, representações vetoriais de palavras, frases ou documentos que levam em conta o contexto em que os termos aparecem, gerados por modelos de linguagem como as Representações de Codificador Bidirecional de Transformadores (BERT, do inglês \textit{Bidirectional Encoder Representations from Transformers}) \parencite{grootendorst2022bertopic}. O BERT considera ao mesmo tempo o contexto à esquerda e à direita de cada palavra \parencite{devlin2019bert}. Com variações como o Sentence-BERT, frases e parágrafos inteiros são convertidos em vetores, e textos de sentido semelhante ficam próximos nesse espaço vetorial \parencite{grootendorst2022bertopic}. São esses mesmos \textit{embeddings} que o RAG usa para encontrar, no material do professor, os trechos mais parecidos com a pergunta do estudante.
}
\novoBlocoTexto{
    A proximidade entre dois vetores pode ser medida pela similaridade de cosseno, que compara a direção dos vetores: quanto mais próximo de 1 for o resultado, maior será a similaridade entre as representações, embora isso não garanta que os textos possuam significados equivalentes. \textcite{karpukhin2020dense} definem a semelhança entre a pergunta e cada passagem pelo produto interno dos seus vetores e observam que, para vetores de comprimento unitário, essa medida equivale ao cosseno. Na recuperação densa, apresentada na seção~\ref{sec:rag}, todas as passagens do material são convertidas em vetores e indexadas com antecedência; quando chega uma pergunta, o sistema calcula o vetor dela e devolve as \textit{k} passagens cujos vetores estão mais próximos, sendo k um número pequeno em relação ao total de passagens \parencite{karpukhin2020dense}.
}
\novoBlocoTexto{
    No entanto, a maior parte dos modelos de \textit{embeddings} foi treinada apenas com textos em inglês, o que limita seu uso em outros idiomas \parencite{wang2024multilingual}. Para contornar essa limitação, \textcite{wang2024multilingual} desenvolveram os modelos \textit{Multilingual E5}, pré-treinados com cerca de 1 bilhão de pares de textos em diversos idiomas e depois ajustados com dados rotulados de alta qualidade, e os avaliaram em recuperação de informação em 16 idiomas e em alinhamento de traduções em mais de 100 idiomas. Segundo os autores, esses modelos podem ser usados em tarefas de recuperação de informação e de similaridade semântica em diferentes idiomas \parencite{wang2024multilingual}. 
}
\subsubsection{{\fontsize{12pt}{14pt}\selectfont Classificação textual baseada em LLM}}
\label{sec:classificacao-textual-llm}
\novoBlocoTexto{
    Uma das metodologias aplicadas à classificação de texto é a classificação zero-shot, na qual o modelo classifica um texto sem nenhum exemplo de treino, contando apenas com os rótulos das categorias \parencite{lobo2025democratizing}. Com LLM, isso é feito por meio de um comando (prompt). \textcite{lobo2025democratizing} informaram à LLM seu papel de analista, a descrição da disciplina e a lista de categorias, e pediram que ela devolvesse apenas o número da categoria de cada pergunta. Em um estudo sobre classificação de perguntas educacionais, \textcite{alfaraby2024analysis} utilizaram o ChatGPT para categorizar 4.959 perguntas em dez classes. Os autores verificaram que o melhor resultado veio do prompt que trazia, além da lista de categorias, a descrição de cada uma, e que essa abordagem é especialmente útil quando não há dados de treinamento \parencite{alfaraby2024analysis}.
}
\novoBlocoTexto{
    É nesse ponto que o material estruturado do professor pode fornecer apoio relevante. Se os tópicos da disciplina forem utilizados como lista fechada de categorias, o problema deixa de ser descobrir os assuntos e passa a ser classificar cada dúvida em assuntos que o professor já reconhece. \textcite{qiu2025systematic} adotaram essa estratégia com uma LLM, instruída a rotular cada conversa com o tópico mais relacionado de uma lista de doze tópicos do curso ou com “Outros” quando não houvesse correspondência. A contagem por tópico revelou volume particularmente alto de mensagens em Probabilidade Condicional, o que coincidia com a dificuldade que o professor já percebia em sala \parencite{qiu2025systematic}. O material oferece ainda a descrição de cada tópico, que melhora a classificação quando incluída no prompt \parencite{alfaraby2024analysis}, e, por meio do RAG, os trechos recuperados carregam metadados como capítulo e página \parencite{gao2023retrieval}, que indicam de qual parte do material a dúvida se aproxima. Para que a LLM não invente um tópico inexistente, a classificação pode ser restringida para que a saída siga obrigatoriamente uma estrutura definida \parencite{geng2023grammar}, como a lista de tópicos do plano de ensino docente.
}
\novoBlocoTexto{
    O desempenho da classificação pode ser avaliado por meio da comparação entre os rótulos atribuídos pelo sistema e aqueles estabelecidos em um conjunto de referência. Entre as métricas utilizadas, destaca-se a acurácia, que representa a proporção total de classificações corretas. A precisão corresponde à proporção de textos atribuídos a determinada categoria que realmente pertencem a ela, enquanto a revocação indica a proporção de textos dessa categoria corretamente identificados pelo sistema \parencite{sokolova2009systematic}. Essas duas medidas são combinadas no F1-score, que expressa seu equilíbrio por meio da média harmônica. Além disso, a matriz de confusão permite identificar quais categorias são classificadas corretamente e quais são confundidas pelo sistema.
}
\novoBlocoTexto{
    Contudo, os resultados apresentados anteriormente nesta seção não significam que a classificação feita por LLM seja perfeita. Em 25 tarefas de processamento de linguagem, o \textit{ChatGPT} apresentou, em média, cerca de 25\% menos qualidade que as melhores soluções especializadas, com perdas maiores nas tarefas mais difíceis \parencite{kocon2023chatgpt}. \textcite{sun2023text} identificaram, nos cenários avaliados, desempenho inferior dos modelos utilizados apenas com instruções em comparação a classificadores treinados especificamente para a tarefa, e seu desempenho costuma ser insuficiente para substituir totalmente a anotação humana \parencite{ziems2024can}. Portanto, a classificação automática das dúvidas deve ser apresentada ao professor como subsídio à análise pedagógica, e não como diagnóstico definitivo de dificuldades de aprendizagem.
}
\subsection{{\fontsize{12pt}{14pt}\selectfont Privacidade e proteção de dados}}
\label{sec:privacidade}
\novoBlocoTexto{
    A Lei Geral de Proteção de Dados Pessoais (LGPD), Lei nº 13.709/2018, define dado pessoal como toda ``informação relacionada a pessoa natural identificada ou identificável" \parencite[art.~5º, I]{brasil2018lgpd}. A adoção de sistemas inteligentes no ensino intensifica a coleta e o tratamento de dados pessoais de estudantes \parencite{miranda2026lgpd}, e os registros de interação em plataformas, assim como os dados de desempenho e comportamento, podem revelar aspectos sensíveis da vida escolar \parencite{vidal2026inteligencia}. 
}
\novoBlocoTexto{
    A LGPD ainda limita o tratamento ``ao mínimo necessário para a realização de suas finalidades" \space \parencite[art.~6º, III]{brasil2018lgpd}. Plataformas educacionais que capturam dados além do necessário para a finalidade pedagógica afrontam esse princípio e ampliam o risco de exposição \parencite{miranda2026lgpd}. Para resolução do problema desta pesquisa, o que mais interessa é saber quais partes do material geram mais dúvidas, e não necessariamente quem as enviou. Essa finalidade pode ser cumprida sem que a identidade do estudante seja exposta. Portanto, a apresentação das informações ao professor é realizada sem identificadores diretos dos estudantes.
}
\novoBlocoTexto{
    A análise do destino para onde as dúvidas discentes são enviadas também deve ser considerada importante. \textcite{vidal2026inteligencia} recomendam verificar, antes de adotar uma ferramenta, o local de armazenamento, o compartilhamento com terceiros e a possibilidade de uso dos dados para treinar modelos. Neste trabalho, o modelo de linguagem é executado localmente, evitando o envio das dúvidas a serviços externos de inferência durante seu processamento. Essa escolha pode reduzir a exposição das informações a terceiros, embora não elimine os riscos associados ao tratamento de dados pessoais, especialmente porque o armazenamento dos registros e dos materiais didáticos utiliza serviços externos. Dessa forma, a execução local da LLM constitui uma medida complementar de proteção, sem dispensar mecanismos de controle de acesso e segurança.
}
\novoBlocoTexto{
    Por fim, nenhuma dessas medidas substitui a transparência. As instituições devem informar aos estudantes para que seus dados serão usados, quem terá acesso a eles e como sua identidade será protegida \parencite{slade2013learning}. No assistente proposto, isso significa informar ao estudante que, ao utilizar a plataforma para fazer uma pergunta, o professor verá um resumo da dúvida e a pergunta original, ambos sem a sua identificação, classificados por tópico da disciplina, com a finalidade de gerar subsídios que possam contribuir para a prática docente. Em síntese, a ocultação da identidade dos estudantes desempenha um papel relevante nesta pesquisa para favorecer a manifestação de dúvidas acadêmicas. Entretanto, a apresentação de informações desidentificadas não elimina completamente o risco de identificação indireta, especialmente em turmas pequenas ou diante de questionamentos muito específicos.
}